\section{Open Challenges and Future Opportunities}
\label{sec:future_directions}

Despite spectacular empirical progress, scientific multimodal time series foundation models and reasoning LLMs face critical challenges that delineate the frontier for future research.

\subsection{Universal Cross-Physics Scientific Tokenizers}
A major architectural barrier is the lack of a universal tokenizer for heterogeneous scientific modalities. Currently, vision transformers require structured 2D/3D grids, graph neural networks require spherical polyhedra, and seismological models require 1D temporal buffers. Although multi-task architectures like UniTS~\cite{gao2024units} and MOMENT~\cite{goswami2024moment} provide foundational steps toward cross-domain time series modeling, developing an omni-modal scientific tokenizer capable of seamlessly mapping sparse in-situ stations, dense gridded reanalysis tensors, asynchronous multi-spectral satellite swaths, continuous acoustic waveforms, and natural language scientific reports into a shared continuous embedding space is a crucial unsolved milestone.

\subsection{Hard Physical Conservation and Extreme Tail Extrapolation}
Most existing foundation models rely on empirical soft loss penalties. When rolled out over extended multi-month or multi-year simulations (e.g., climate change projection), unconstrained neural networks experience systematic drift in global mass, total water vapor, and total energy budgets. Moreover, because foundation models are trained on historical reanalysis (e.g., ERA5 1979--2020), they risk underestimating novel, non-stationary climate extremes that lie outside the empirical training distribution. While generative diffusion models like GenCast~\cite{price2023gencast} and SEEDS~\cite{li2024seeds} significantly improve extreme-event tail probabilities, integrating hard mathematical conservation operators (e.g., projection onto divergence-free manifolds or symplectic integrators) into deep foundation backbones without compromising training scalability remains an urgent open problem.

\subsection{Hybrid Numerical-Neural Planetary Digital Twins}
Rather than viewing deep learning and numerical differential equation solvers as competitors, the next decade will likely be defined by hybrid numerical-neural architectures (exemplified by NeuralGCM~\cite{kochkov2024neuralgcm} and the ClimSim benchmark~\cite{yu2023climsim}). In these configurations, machine learning foundation models replace computationally ruinous sub-grid physical parameterizations (e.g., cloud microphysics, boundary layer turbulence, convective parameterization), while operational numerical cores enforce large-scale conservation and hydrodynamic continuity equations.

\subsection{Trustworthy Autonomous Scientific Discovery}
While autonomous multi-agent frameworks (such as ClimateAgent~\cite{li2025climateagent} and OceanGPT~\cite{bi2024oceangpt}) demonstrate immense potential, deploying LLM agents in operational scientific decision-making introduces risks of physical hallucination. Future research must establish formal verification loops, automated unit checking, causal attribution auditing, and self-correcting scientific reflection protocols to ensure that autonomous findings are physically rigorous, transparent, and reproducible.
